\chapter{Producing a submission}\label{ch:submission}

\begin{onething}
An evaluation result is a committed \code{anomaly_detection_prep} exported as a boolean 0/1 GeoTIFF plus a point
shapefile, with no score anywhere a judge could read as a heatmap, on the product's own map grid. For PRISMA that grid
now comes from the HE5's UTM corner attributes read as pixel \emph{centres}, pitch $=$ span$/(n-1)$ and origin $=$
upper-left centre minus half a pixel, replacing a latitude and longitude bounding box that was a median 2.4 pixels
off.
\end{onething}

On 2026-10-31 an AIGC Stage-III evaluation circular set the rules for the deliverables, with a report due on
2~November 2026 and a presentation on 4~November. The plan and its status are
\file{docs/proposals/aigc-evaluation-submission.md}: its changes C1--C6 were built and tested that day, and C7, the
rehearsal, was done on the \emph{local} stack (``images rebuilt, worker \code{running}, export rehearsed on a committed
PRISMA prep: \code{hyper_*} bundle, uint8 \{0,1\} mask in the product's UTM zone''). A production or air-gapped
deployment is not approved and is a separate step. This chapter is how the product meets the rules, and the one piece
of geometry that had to change.

\begin{analogy}
A land registry accepts boundary maps only as surveyed lines on its own grid: not a sketch, and not a shaded map of
``likely boundaries''. You may file up to three claims per plot. Until this week the office drew PRISMA plots by eye
from the corners of a road atlas; now it reads the surveyor's grid references off the deed.
\end{analogy}

%======================================================================
\section{The rules, and the one bundle that meets them}\label{sec:submission:rules}

\begin{center}\small
\begin{tabular}{@{}lp{0.38\linewidth}p{0.44\linewidth}@{}}
\toprule
& rule (circular) & how the product meets it\\
\midrule
R2a & boolean 0/1 anomaly map as a GeoTIFF with a real CRS & the prep commit writes uint8 0/1; the export stamps on
the raw scene's georeference\\
R2b & the same map as a shapefile & one point per anomaly pixel, full sidecar set (.shp .shx .dbf .prj .cpg)\\
R2c & no heatmap-style output: it disqualifies & no score raster in the bundle, and since 2026-10-31 no score
attribute in the shapefile\\
R2d & up to three results per scene & three committed prep actions, three exports\\
R2e & library matching by a foundation model and by classic SAM, compared in the report & the per-object table,
\cref{sec:submission:objects}\\
R1 & a reflectance product from L1 data & not in the product; a writer for another module's output,
\cref{sec:submission:reflectance}\\
\bottomrule
\end{tabular}
\end{center}

Only the prep bundle qualifies: the export of a committed \code{anomaly_detection_prep}, built by
\code{build_thermal_bundle} in \file{app/anomaly_detection/export.py}. The project export zips every output tree,
\file{anomaly_score.tif} (a heatmap) included; the materials bundle of \cref{sec:material:export} carries int16 family
codes, not 0/1. Two changes were made to the prep bundle (C2):
\begin{lstlisting}
prefix = "thermal" if spec.sensor_type in THERMAL_SENSORS else "hyper"   # ("landsat9", "hotsat1")
folder = f"{prefix}_{_short(spec.scene_id)}_{_short(spec.action_id)}"
keep_cols = ["row", "col", "latitude", "longitude", "geometry"]           # "score" removed
\end{lstlisting}
A PRISMA bundle used to be named \code{thermal_*}, against the rule that a name says \code{hyper} or \code{thermal}.
The comment on the dropped column: ``No score column: the submission map is boolean, and a continuous score reads as a
heatmap (AIGC rules, 2026-10-31). The CSV keeps it.''

\begin{sourcenote}[: the score is still in the zip]
The CSV, \code{<prefix>_anomaly_summary.csv}, is inside the same zip with its \code{score} column; the proposal calls it
``not submitted'', so leave it out of the upload. Two names were not updated: the function is still
\code{build_thermal_bundle}, and the manifest's schema is still \code{allotrope-export-thermal/1} for a \code{hyper_}
bundle.
\end{sourcenote}

%======================================================================
\section{The PRISMA grid}\label{sec:submission:grid}

Recall \cref{eq:material:affine}: $X=a\,j+b\,i+c$, $Y=d\,j+e\,i+f$, with $(j,i)$ a position in pixel units measured from
the top-left \emph{corner} of the image, so the top-left pixel's centre is $(0.5,0.5)$. Until 2026-10-31 a PRISMA export
used the latitude and longitude bounding box of \cref{we:material:georef}. The ground-truth rasters are in UTM. On Set-1
the box put pixels a median 2.4\,px from the product grid, 95\% within 4.1, at most 5.6
\tagM{docs/proposals/aigc-evaluation-submission.md}: at 30\,m, $2.4\times30=72$\,m typical and $5.6\times30=168$\,m at
worst. The internal evaluation never noticed, because it compares masks in pixel space.

A PRISMA L2D product is already gridded in UTM, and the HE5's root attributes say how (\code{prisma_utm_grid},
\file{app/georef/from_scene_file.py}): \code{Epsg_Code}, and the easting and northing of the upper-left, upper-right and
lower-left corner pixels. Its docstring: they ``are the corner pixels' \emph{centres}, so the pitch is the corner span
over $(n-1)$ and the origin is the UL centre minus half a pixel''.

\paragraph{Centres, not edges.} Five fence posts 30\,m apart span $4\times30=120$\,m from the first post's centre to the
last's. In general, $n$ pixel centres are $n-1$ pitches apart, so with $W$ columns
\begin{align*}
\text{UR}_E-\text{UL}_E &= (W-1)\,\Delta x &&\text{centres of the first and last column}\\
\Delta x &= \frac{\text{UR}_E-\text{UL}_E}{W-1} &&\text{divide by } W-1\\
c &= \text{UL}_E-\tfrac12\Delta x &&\text{the image edge is half a pixel left of the first centre,}
\end{align*}
and the same for rows with the northings, except that rows run southwards, so $e=-\Delta y$ and
$f=\text{UL}_N+\tfrac12\Delta y$.

\begin{worked}[label={we:submission:set1}]{Set-1's grid by hand}
The test file \file{tests/test_utils/test_georef_prisma.py} holds Set-1's attributes, stored as 32-bit floats like the
file: \code{Epsg_Code} 32642; UL easting 545,022.75 and northing 2,501,422.5; UR easting 580,872.5; LL northing
2,465,872.5. The image is $H=1186$ rows by $W=1196$ columns.
\begin{align*}
\text{east span} &= 580{,}872.5-545{,}022.75 = 35{,}849.75\text{ m},\\
\Delta x &= 35{,}849.75/1195 = 29.99979 \ \to\ 30.0\text{ m (rounded to centimetres)},\\
\text{north span} &= 2{,}501{,}422.5-2{,}465{,}872.5 = 35{,}550\text{ m},\\
\Delta y &= 35{,}550/1185 = 30.0\text{ m},\\
c &= 545{,}022.75-15 = 545{,}007.75,\qquad f = 2{,}501{,}422.5+15 = 2{,}501{,}437.5 .
\end{align*}
The transform is $(30,\ 0,\ 545{,}007.75,\ 0,\ -30,\ 2{,}501{,}437.5)$ in EPSG:32642 (UTM zone 42 north): exactly the
ground-truth raster's grid, which the test asserts. Check the top-left centre: $X=30\times0.5+545{,}007.75=545{,}022.75$,
the UL attribute \tagC.
\end{worked}

Read the attributes as edges instead and $\Delta x=35{,}849.75/1196=29.9747$\,m: 0.025\,m short per pixel, which over
1196 columns adds up to $1196\times0.0253\approx30$\,m, one whole pixel at the far edge. The rounding to centimetres
(the code's comment: ``float32 attrs: 29.9998 -> 30.0'') snaps the pitch to the product's 30\,m; its price is that the
bottom-right centre comes out at 580,872.75 instead of the attribute's 580,872.5, a quarter of a metre, under 1\% of a
pixel \tagC.

\paragraph{When the old method still runs.} \code{_resolve_prisma} uses the product grid only when the attributes are
present and sane (a positive EPSG code, at least $2\times2$ pixels, positive spans) and the vendable's shape equals the
latitude array's. Otherwise it falls back to the bounding box, whose docstring now says ``approximate, can be pixels
off''; a test checks both refusals. The same function serves the reflectance writer's \code{--georef-from}.

%======================================================================
\section{Native bands by default}\label{sec:submission:native}

The proposal's gap~3 was ``Indradhanu V2 cannot be run from the UI, because no UI control sets the band filter to native
bands''. Its fix, C3: every hyperspectral sensor default in \file{band_filter_apply.py} now reads
\begin{lstlisting}
"use_common_wavelength_grid": False,   # native bands (2026-10-31): V2 reads them; 165-grid models resample at scoring
\end{lstlisting}
The dialog prefills its form from these defaults, so the default is now V2's setting, in line with the standing design
decision ``One inference method: native bands'' (\file{docs/lld/spectral-stem.md}, 2026-10-15). The old 165-band model
keeps working because \code{anomaly_scoring} resamples for it, per model, with the band filter's own resampler
(\code{_to_common_grid}; a cube already on the grid is passed through):

\begin{center}\small
\begin{tabular}{@{}lccc@{}}
\toprule
band filter & v0.2.0 (165 grid) & V2 (native) & MNF-RX\\
\midrule
native (the default) & runs, resampled & runs & runs\\
165 grid & runs & refused & runs\\
\bottomrule
\end{tabular}
\end{center}

So one scoring action can now run both Indradhanus on a native band filter (\cref{sec:material:v2}).

\begin{sourcenote}[: a field default still describes the 165-band world]
The pydantic field \code{BandFilterApplyConfig.use_common_wavelength_grid} still has \code{default=True}, and its
description still says the grid is ``Required by Indradhanu (v0.2.0)''. The UI sends the sensor defaults, so the UI gets
native bands; \textbf{an API call that omits the field gets the 165 grid, and V2 then fails}
(\file{docs/09-known-issues.md} records it as deliberate). The worker's fallbacks \code{cfg.get("use_common_wavelength_grid", True)}
agree with the field, not with the defaults. Two more comments were stale when this chapter was first written,
\code{_check_band_grid}'s docstring and the resolver's comment on \code{band_grid}; both were corrected on
2026-10-31 and now describe the refusal of V2 on the 165 grid and the resampling for a common165 model.
\end{sourcenote}

%======================================================================
\section{Three results per scene}\label{sec:submission:three}

Each result is one committed prep action, so three results are a procedure, not code: three prep actions (they may
share one scoring output), three commits, three exports, each zip named by its action's short id. A re-commit overwrites
a prep's mask, so one prep cannot carry two results. The proposal suggests, per hyperspectral scene and all with
residual Mahalanobis: Indradhanu~V2 (Heavy); V2 and MNF-RX combined at the prep weights; and MNF-RX alone or V2 at a
second threshold. Choosing the thresholds stays a person's job (\cref{sec:material:why}); with no ground truth on the
evaluation scenes, the proposal asks whether to fix a percentile per model from the four labelled PRISMA scenes.

%======================================================================
\section{The per-object table}\label{sec:submission:objects}

The report must compare a foundation model with SAM, per object. \file{scripts/aigc_object_table.py} groups the
committed mask into connected components with \code{scipy.ndimage.label} and a $3\times3$ structure of ones.

\begin{definition}[4- and 8-connectivity]
Two flagged pixels are \emph{4-connected}\index{connectivity} if they share an edge, and \emph{8-connected} if they share
an edge or a corner. An object is a largest set of flagged pixels linked by a chain of such neighbours.
\end{definition}

\begin{worked}[label={we:submission:components}]{Two objects or three}
\[
\begin{pmatrix}1&1&0&0&0\\0&1&0&0&1\\0&0&1&0&1\\0&0&0&0&0\end{pmatrix}
\]
The pixel at row 2, column 2 touches row 1, column 1 only at a corner. With 8-connectivity it joins the top-left group:
2 objects. With 4-connectivity it is alone: 3 objects. \code{label} agrees in both cases \tagC.
\end{worked}

Per object the script writes size, centroid (latitude and longitude through the scene's own georeference, so PRISMA uses
\cref{sec:submission:grid}'s grid), and three opinions: SAM (\file{matches.parquet}, rank~0: the majority top-1 name and
its share, its chapter, the median angle, the three commonest names), the family match (\file{family_matches.parquet},
answered pixels: majority family, group, mean confidence, answered share) and, with \code{--namer-ckpt}, the namer of
\cref{ch:namer} (mean $P(M)$, majority class and share, mean $P(k\mid M)$, mean fill). Outputs: \file{objects.csv},
\file{objects.geojson}, and a \file{findings.md} stub with agreement counts. Its docstring is honest that the family
match is ``calibrated Simplex ACE (classical)'': the namer is the only foundation-model column, read from an N001
checkpoint (\file{ep<n>.pt} exists after every epoch, so a preliminary one can be used before the block ends).

%======================================================================
\section{The reflectance writer}\label{sec:submission:reflectance}

R1 asks for surface reflectance from L1 data. The atmospheric correction is another session's module;
\file{scripts/write_reflectance_product.py} only writes its result: GeoTIFF and/or ENVI BSQ with a \file{.hdr}, every
wavelength recorded, water-absorption bands removed (default 1340--1460 and 1790--1960\,nm, ends included), and
reflectance $\times10000$ as int16 with nodata $-9999$:
\begin{lstlisting}
data = np.where(np.isfinite(cube), np.clip(np.round(cube * SCALE), -9998, 32767), NODATA).astype(np.int16)
\end{lstlisting}

\begin{worked}[label={we:submission:scale}]{Writing five values}
\begin{center}\small
\begin{tabular}{@{}rrrl@{}}
\toprule
reflectance & $\times10000$ & written & why\\
\midrule
0.2 & 2000 & 2000 & \\
0.12345 & 1234.5 & 1234 & \code{np.round} sends a half to the even neighbour\\
0.12355 & 1235.5 & 1236 & here the even neighbour is above\\
NaN & NaN & $-9999$ & nodata\\
4.0 & 40000 & 32767 & clipped to int16's top\\
\bottomrule
\end{tabular}
\end{center}
The lower clip is $-9998$, so no real value collides with nodata. Bands at 1330, 1345, 1460 and 1470\,nm under the
default drop: 1345 and 1460 go (1460 is an end of the range), 1330 and 1470 stay \tagC.
\end{worked}

\begin{drift}
The submission's status (C1--C7 on the local stack, deployment not approved), the deadlines, the Set-1 offset
(median 2.4\,px, 95\% 4.1, max 5.6), the three suggested results, and the field default and its description in
\cref{sec:submission:native}, which a later change may align.
\end{drift}

%======================================================================
\begin{confusions}
\confusion{Any export from the product will do.}{Only the prep bundle. The project export carries a heatmap; the
materials raster is family codes.}
\confusion{No score in the shapefile means no score in the zip.}{The CSV in the same zip keeps it. Do not upload it.}
\confusion{PRISMA's corner attributes are the image's corners.}{They are the corner pixels' centres. Read as edges, the
far edge drifts by a pixel.}
\confusion{The georeference error moved the benchmark.}{The evaluation compares masks in pixel space. Only deliverables
overlaid on a map were affected.}
\confusion{V2 still needs a band-filter setting.}{Not from the UI since 2026-10-31. An API call that omits the field
still gets the 165 grid.}
\confusion{The reflectance script corrects the atmosphere.}{It writes another module's output. ``The atmospheric
correction itself is not here.''}
\end{confusions}

\begin{checkyourself}
\item\label{ex:submission:1} A PRISMA product has $W=1000$ columns, UL easting 400,015 and UR easting 429,985. Compute
  $\Delta x$ and $c$. What pitch does dividing by $W$ give, and how far off is the last column's centre then?
\item\label{ex:submission:2} With Set-1's transform, where is the centre of the pixel in row 10, column 20?
\item\label{ex:submission:3} A band filter was created through the API without \code{use_common_wavelength_grid}. A
  scoring action picks \code{["Indradhanu V2", "Indradhanu"]}. What happens?
\item\label{ex:submission:4} Count the objects in $\begin{pmatrix}1&0&1\\0&1&0\\1&0&1\end{pmatrix}$ under 8- and under
  4-connectivity.
\item\label{ex:submission:5} Write 0.03, $-1.2$ and 3.5 as the reflectance writer would.
\item\label{ex:submission:6} Why is a second commit on the same prep action not a way to make result~2?
\end{checkyourself}

\par\bigskip\noindent\textit{That brings the book to 2026-10-31: Indradhanu~V2 (Heavy) in the product, native bands by
default, a spectral namer in training, and the first evaluation submission on the product's own grid.}
